The standard environment for Hope, implicitly used by every session
and module.
\subsection{Standard type constructors}

Standard type variables.
\begin{verbatim}
    typevar alpha, beta, gamma;
\end{verbatim}
Function and product types.
\begin{verbatim}
    infixr -> : 2;
    abstype neg -> pos;

    infixr # : 4;
    abstype pos # pos;

    --- (a # b) (x, y) <= (a x, b y);

    infixr X : 4;
    type alpha X beta == alpha # beta;
\end{verbatim}
Normal right-to-left composition of f and g.
\begin{verbatim}
    infix o : 2;
    dec o : (beta -> gamma) # (alpha -> beta) -> alpha -> gamma;
    --- (f o g) x <= f(g x);

    --- (a -> b) f <= b o f o a;
\end{verbatim}
The identity function.
\begin{verbatim}
    dec id : alpha -> alpha;
    --- id x <= x;
\end{verbatim}
Booleans with the standard operations.
\begin{verbatim}
    data bool == false ++ true;
    type truval == bool;

    infix or : 1;
    infix and : 2;

    dec not : bool -> bool;
    --- not true <= false;
    --- not false <= true;

    dec and : bool # bool -> bool;
    --- false and p <= false;
    --- true and p <= p;

    dec or : bool # bool -> bool;
    --- true or p <= true;
    --- false or p <= p;

    dec if_then_else : bool -> alpha -> alpha -> alpha;
    --- if true then x else y <= x;
    --- if false then x else y <= y;
\end{verbatim}
Lists.
\begin{verbatim}
    infixr :: : 5;
    data list alpha == nil ++ alpha :: list alpha;
\end{verbatim}
List concatenation.
\begin{verbatim}
    infixr <> : 5;
    dec <> : list alpha # list alpha -> list alpha;
    --- [] <> ys <= ys;
    --- (x::xs) <> ys <= x::(xs <> ys);
\end{verbatim}
The type num behaves as if it were defined by:
\begin{verbatim}
    !        data num == 0 ++ succ num;
\end{verbatim}
but is actually implemented a little more efficiently.
The type also contains negative numbers (and reals in some implementations).
\begin{verbatim}
    data num == succ num;
\end{verbatim}
Similarly, the type char behaves as if it were defined as an enumeration
(in the normal order) of all the ASCII character constants.
\begin{verbatim}
    abstype char;

    --- char x <= x;
\end{verbatim}
\subsection{Internally defined functions}

Comparison functions:
these evaluate their arguments only to the extent necessary to determine
their order.  In comparing distinct constants and constructors, the
lesser is the one which came earlier in the data definition in which
both were defined.  Pairs compare lexicographically, while comparison
of functions triggers a run-time error.  On the types num, char and
list(char), these rules give the normal orderings.
\begin{verbatim}
    infix =, /= : 3;
    infix <, =<, >, >= : 4;
    dec =, /=, <, =<, >, >= : alpha # alpha -> bool;
\end{verbatim}
Lower-level comparison functions.
\begin{verbatim}
    data relation == LESS ++ EQUAL ++ GREATER;

    dec compare : alpha # alpha -> relation;

    --- x =  y <= (\ EQUAL => true | _ => false) (compare(x, y));
    --- x /= y <= (\ EQUAL => false | _ => true) (compare(x, y));
    --- x <  y <= (\ LESS => true | _ => false) (compare(x, y));
    --- x >= y <= (\ LESS => false | _ => true) (compare(x, y));
    --- x >  y <= (\ GREATER => true | _ => false) (compare(x, y));
    --- x =< y <= (\ GREATER => false | _ => true) (compare(x, y));
\end{verbatim}
Conversions.
\begin{verbatim}
    dec ord : char -> num;
    dec chr : num -> char;

    dec num2str : num -> list char;
    dec str2num : list char -> num;
\end{verbatim}
The contents of a named file (created lazily).
\begin{verbatim}
    dec read : list char -> list char;
\end{verbatim}
Abort execution with an error message.
\begin{verbatim}
    dec error : list char -> alpha;
\end{verbatim}
The usual arithmetical functions.
\begin{verbatim}
    infix +, - : 5;
    infix *, /, div, mod : 6;
    dec +, -, *, /, div, mod : num # num -> num;
\end{verbatim}
Math library.
\begin{verbatim}
    dec cos, sin, tan, acos, asin, atan : num -> num;
    dec atan2, hypot : num # num -> num;
    dec cosh, sinh, tanh, acosh, asinh, atanh : num -> num;
    dec abs, ceil, floor : num -> num;
    dec exp, log, log10, sqrt : num -> num;
    dec pow : num # num -> num;
    dec erf, erfc : num -> num;
\end{verbatim}
Any extra arguments on the interpreter command line.
\begin{verbatim}
    dec argv : list(list char);
\end{verbatim}
Part of the special treatment of the succ constructor.
\begin{verbatim}
    dec succ : num -> num;
    --- succ n <= n+1;
\end{verbatim}
